\documentclass[10pt,a4paper]{amsart}
\usepackage{amsmath,amssymb,amsthm,graphicx}
% Retained minimal preamble for the Stacks Project chapter varieties.tex.
% Verbatim excerpt of the chapter preamble: Xy-pic setup (native preamble lines 14-18), the native theorem
% declarations with the mathematical macros (native preamble lines 154-217), and the project own no-op
% definitions of the reference/slogan/history display environments (native wrapper.tex lines 1-3) that the
% counted statement uses for its slogan. Copied mechanically; nothing added, removed or redefined.
% For commutative diagrams we use Xy-pic
\usepackage[all]{xy}

% We use 2cell for 2-commutative diagrams.
\xyoption{2cell}
% Theorem environments.
%
\theoremstyle{plain}
\newtheorem{theorem}[subsection]{Theorem}
\newtheorem{proposition}[subsection]{Proposition}
\newtheorem{lemma}[subsection]{Lemma}

\theoremstyle{definition}
\newtheorem{definition}[subsection]{Definition}
\newtheorem{example}[subsection]{Example}
\newtheorem{exercise}[subsection]{Exercise}
\newtheorem{situation}[subsection]{Situation}

\theoremstyle{remark}
\newtheorem{remark}[subsection]{Remark}
\newtheorem{remarks}[subsection]{Remarks}

\numberwithin{equation}{subsection}

% Macros
%
\def\lim{\mathop{\mathrm{lim}}\nolimits}
\def\colim{\mathop{\mathrm{colim}}\nolimits}
\def\Spec{\mathop{\mathrm{Spec}}}
\def\Hom{\mathop{\mathrm{Hom}}\nolimits}
\def\Ext{\mathop{\mathrm{Ext}}\nolimits}
\def\SheafHom{\mathop{\mathcal{H}\!\mathit{om}}\nolimits}
\def\SheafExt{\mathop{\mathcal{E}\!\mathit{xt}}\nolimits}
\def\Sch{\mathit{Sch}}
\def\Mor{\mathop{\mathrm{Mor}}\nolimits}
\def\Ob{\mathop{\mathrm{Ob}}\nolimits}
\def\Sh{\mathop{\mathit{Sh}}\nolimits}
\def\NL{\mathop{N\!L}\nolimits}
\def\CH{\mathop{\mathrm{CH}}\nolimits}
\def\proetale{{pro\text{-}\acute{e}tale}}
\def\etale{{\acute{e}tale}}
\def\QCoh{\mathit{QCoh}}
\def\Ker{\mathop{\mathrm{Ker}}}
\def\Im{\mathop{\mathrm{Im}}}
\def\Coker{\mathop{\mathrm{Coker}}}
\def\Coim{\mathop{\mathrm{Coim}}}

% Boxtimes
%
\DeclareMathSymbol{\boxtimes}{\mathbin}{AMSa}{"02}

%
% Macros for moduli stacks/spaces
%
\def\QCohstack{\mathcal{QC}\!\mathit{oh}}
\def\Cohstack{\mathcal{C}\!\mathit{oh}}
\def\Spacesstack{\mathcal{S}\!\mathit{paces}}
\def\Quotfunctor{\mathrm{Quot}}
\def\Hilbfunctor{\mathrm{Hilb}}
\def\Curvesstack{\mathcal{C}\!\mathit{urves}}
\def\Polarizedstack{\mathcal{P}\!\mathit{olarized}}
\def\Complexesstack{\mathcal{C}\!\mathit{omplexes}}
% \Pic is the operator that assigns to X its picard group, usage \Pic(X)
% \Picardstack_{X/B} denotes the Picard stack of X over B
% \Picardfunctor_{X/B} denotes the Picard functor of X over B
\def\Pic{\mathop{\mathrm{Pic}}\nolimits}
\def\Picardstack{\mathcal{P}\!\mathit{ic}}
\def\Picardfunctor{\mathrm{Pic}}
\def\Deformationcategory{\mathcal{D}\!\mathit{ef}}
\newenvironment{reference}{}{}
\newenvironment{slogan}{}{}
\newenvironment{history}{}{}
\usepackage{cleveref}
\crefname{originallemma}{original lemma}{original lemmas}
\Crefname{originallemma}{Original lemma}{Original lemmas}
\makeatletter
\newlabel{lemma-dim-1-proper-projective}{{0A26}{source}{}{sourcepointer-1}{}}
\newlabel{lemma-dim-1-proper-projective@cref}{{[originallemma][0][]0A26}{[1][0][]source}{}{}{}}
\makeatother
\crefname{originalreference}{original reference}{original references}
\Crefname{originalreference}{Original reference}{Original references}
\makeatletter
\newlabel{coherent-lemma-vanshing-gives-ample}{{0B5U}{source}{}{sourcepointer-2}{}}
\newlabel{coherent-lemma-vanshing-gives-ample@cref}{{[originalreference][0][]0B5U}{[1][0][]source}{}{}{}}
\makeatother
\crefname{originallemma}{original lemma}{original lemmas}
\Crefname{originallemma}{Original lemma}{Original lemmas}
\makeatletter
\newlabel{lemma-dense-smooth-open-variety-over-perfect-field}{{0B8X}{source}{}{sourcepointer-3}{}}
\newlabel{lemma-dense-smooth-open-variety-over-perfect-field@cref}{{[originallemma][0][]0B8X}{[1][0][]source}{}{}{}}
\makeatother
\crefname{originallemma}{original lemma}{original lemmas}
\Crefname{originallemma}{Original lemma}{Original lemmas}
\makeatletter
\newlabel{lemma-regular-point-on-curve}{{0B8Y}{source}{}{sourcepointer-4}{}}
\newlabel{lemma-regular-point-on-curve@cref}{{[originallemma][0][]0B8Y}{[1][0][]source}{}{}{}}
\makeatother
\title{Existence of invertible sheaves on a proper curve with prescribed vanishing of H\textasciicircum{}0 and H\textasciicircum{}1}
\author{The Stacks Project Authors}
\date{stacks-project:varieties.tex:local-snapshot-20261011; GNU FDL 1.2 or later}
\begin{document}
\maketitle
\noindent\textit{Editorial note.} All delivered mathematical text is the author's own, copied verbatim from the pinned Stacks Project chapter file varieties.tex; no statement or proof step was rewritten, reordered or removed.\par\medskip
\noindent\textit{Editorial note.} The counted span cites four cross-chapter labels, each recorded in assembly.external\_references with its canonical tag as a navigation pointer.\par\medskip
\noindent\textit{Editorial note.} The proof cites no bibliography item inside the counted span.\par\medskip
\noindent\textit{Allowed standard prerequisites (editorial).}\begin{itemize}
\item Existence of a nonempty open set of nonsingular points. Stacks Project, Varieties, Lemma lemma-dense-smooth-open-variety-over-perfect-field (tag 0B2J). Supplies the open U used to choose the point x.
\item The regular-point criterion and the ideal sheaf of a closed point of a curve. Stacks Project, Varieties, Lemma lemma-regular-point-on-curve (tag 0B2S). Identifies the ideal sheaf of x as invertible, which makes the exact sequence computation work.
\end{itemize}\medskip
\ifdefined\hypertarget\hypertarget{sourcepointer-1}{}\fi
\noindent\textit{Source reference.} Original lemma 0A26: sources/stacks-chow/tags.txt, source lines 12694--12694. Cross-chapter standard prerequisite cited by the author inside the counted span; it is used as a pointer and its own proof is not part of this item.\par\medskip
\ifdefined\hypertarget\hypertarget{sourcepointer-2}{}\fi
\noindent\textit{Source reference.} Original reference 0B5U: sources/stacks-chow/tags.txt, source lines 14047--14047. Cross-chapter standard prerequisite cited by the author inside the counted span; it is used as a pointer and its own proof is not part of this item.\par\medskip
\ifdefined\hypertarget\hypertarget{sourcepointer-3}{}\fi
\noindent\textit{Source reference.} Original lemma 0B8X: sources/stacks-chow/tags.txt, source lines 14155--14155. Cross-chapter standard prerequisite cited by the author inside the counted span; it is used as a pointer and its own proof is not part of this item.\par\medskip
\ifdefined\hypertarget\hypertarget{sourcepointer-4}{}\fi
\noindent\textit{Source reference.} Original lemma 0B8Y: sources/stacks-chow/tags.txt, source lines 14156--14156. Cross-chapter standard prerequisite cited by the author inside the counted span; it is used as a pointer and its own proof is not part of this item.\par\medskip
\begin{lemma}
\label{lemma-general-degree-g-line-bundle}
Let $k$ be an algebraically closed field. Let $X$ be a proper curve over $k$.
Then there exist
\begin{enumerate}
\item an invertible $\mathcal{O}_X$-module $\mathcal{L}$ with
$\dim_k H^0(X, \mathcal{L}) = 1$ and $H^1(X, \mathcal{L}) = 0$, and
\item an invertible $\mathcal{O}_X$-module $\mathcal{N}$ with
$\dim_k H^0(X, \mathcal{N}) = 0$ and $H^1(X, \mathcal{N}) = 0$.
\end{enumerate}
\end{lemma}

\begin{proof}
Choose a closed immersion $i : X \to \mathbf{P}^n_k$
(Lemma \ref{lemma-dim-1-proper-projective}).
Setting $\mathcal{L} = i^*\mathcal{O}_{\mathbf{P}^n}(d)$
for $d \gg 0$ we see that there exists an invertible sheaf
$\mathcal{L}$ with $H^0(X, \mathcal{L}) \not = 0$ and
$H^1(X, \mathcal{L}) = 0$ (see
Cohomology of Schemes, Lemma \ref{coherent-lemma-vanshing-gives-ample}
for vanishing and the references therein for nonvanishing).
We will finish the proof of (1) by descending induction on
$t = \dim_k H^0(X, \mathcal{L})$. The base case $t = 1$ is trivial.
Assume $t > 1$.

\medskip\noindent
Let $U \subset X$ be the nonempty open subset of nonsingular points
studied in Lemma \ref{lemma-dense-smooth-open-variety-over-perfect-field}.
Let $s \in H^0(X, \mathcal{L})$ be nonzero. There exists a closed
point $x \in U$ such that $s$ does not vanish in $x$. Let $\mathcal{I}$
be the ideal sheaf of $i : x \to X$ as in
Lemma \ref{lemma-regular-point-on-curve}. Look at the
short exact sequence
$$
0 \to \mathcal{I} \otimes_{\mathcal{O}_X} \mathcal{L} \to
\mathcal{L} \to i_*i^*\mathcal{L} \to 0
$$
Observe that $H^0(X, i_*i^*\mathcal{L}) = H^0(x, i^*\mathcal{L})$
has dimension $1$ as $x$ is a $k$-rational point ($k$ is algebraically
closed). Since $s$ does not vanish at $x$ we conclude that
$$
H^0(X, \mathcal{L}) \longrightarrow H^0(X, i_*i^*\mathcal{L})
$$
is surjective. Hence
$\dim_k H^0(X, \mathcal{I} \otimes_{\mathcal{O}_X} \mathcal{L}) = t - 1$.
Finally, the long exact sequence of cohomology also shows that
$H^1(X, \mathcal{I} \otimes_{\mathcal{O}_X} \mathcal{L}) = 0$
thereby finishing the proof of the induction step.

\medskip\noindent
To get an invertible sheaf as in (2) take an invertible sheaf
$\mathcal{L}$ as in (1) and do the argument in the previous paragraph
one more time.
\end{proof}

\end{document}
